\reportsection{Definitions}
\label{sec:approach}

\subsection{The Object and Its Merkle Instantiation}
\label{sec:syntax}

The mathematical boundary should not drift from the reference the
implementation follows, so we use the Chiesa--Yogev notation and algorithm
split~\cite{CY26}.  We fix the abstract object first and immediately recover the
Merkle commitment as its hash-based realisation; keeping the two at the same
level of abstraction is what lets the implementation claim ``$\textsf{ark-mt}$
is a vector commitment.''  The properties follow in \cref{sec:props}, and
\cref{sec:results} realises both as executable contracts.

\begin{definition}[Vector commitment scheme]\label{def:vcgeneral}
A \emph{vector commitment scheme} over alphabet $\MTAlphabet$ and length
$\MTMessageLength$ is a triple $\VC=(\Commit,\Open,\VCCheck)$ that binds a
message $\msgvec\in\MTAlphabet^{\MTMessageLength}$ to a short string $\comm$ and
authenticates the entries at an index set $\idxset\subseteq[\MTMessageLength]$:
\begin{description}[leftmargin=2.3em,style=nextline]
  \item[$\Commit(\msgvec)\to(\comm,\aux)$.]
    Commit to $\msgvec$, retaining committer state $\aux$ for later openings.
  \item[$\Open(\aux,\idxset)\to\opening$.]
    Authenticate the locations $\idxset\subseteq[\MTMessageLength]$.
  \item[$\VCCheck(\comm,\idxset,\MTMessageSubVector,\opening)\to\bits$.]
    Decide whether $\opening$ authenticates the subvector
    $\MTMessageSubVector\in\MTAlphabet^{\idxset}$ at those locations against
    $\comm$.
\end{description}
A scheme must be \emph{complete} and \emph{position binding}; it is
\emph{optionally hiding}---the zero-knowledge case---when openings must leak
nothing about unopened entries (\cref{sec:props}).
\end{definition}

\begin{definition}[Merkle commitment scheme]\label{def:vc}
The \emph{Merkle commitment scheme}
$\MTSymbol=(\MTCommit,\MTOpen,\MTCheck)$ instantiates \cref{def:vcgeneral} in
the random-oracle model: the three algorithms take query access to
$\ROFunction\in\RODistribution{\ROOutputSize}$, the commitment $\MTCommitment$
is the tree root, and the committer state $\MTTrapdoor$ is the retained salts
and labels.  Two parameters extend the alphabet $\MTAlphabet$ and length
$\MTMessageLength$: the random-oracle output size $\ROOutputSize\in\N$ and the
salt size $\MTSaltSize\in\N$, written
$\MTConstructor{\ROOutputSize}{\MTAlphabet}{\MTMessageLength}{\MTSaltSize}$ when
explicit.  It is \emph{optionally hiding} exactly as in \cref{def:vcgeneral},
with hiding carried by $\MTSaltSize$.
\end{definition}

The Merkle scheme is a \emph{subset} of \cref{def:vcgeneral}: every abstract
algorithm has a Merkle refinement and the properties of \cref{sec:props}
specialise directly, with digest length $\ROOutputSize$ governing binding and
salt cardinality $\MTSaltSize$ governing hiding.  The abstract triple states
\emph{what} a vector commitment guarantees; the pair
$(\ROOutputSize,\MTSaltSize)$ states \emph{how much} each guarantee is worth,
and no tuning of one supplies what the other lacks.

There is no cryptographic $\Setup$: the public parameters are the hash
function, choosing it draws no secret randomness, so the setup is transparent
and the trapdoor is $\td=\bot$.  Two consequences follow.  Extraction is
ROM-only, a query-trace argument that reduces to binding; and with no trapdoor
to equivocate with, a simulator needing standard-model equivocation cannot use
a Merkle commitment.

\subsection{Completeness, Binding, and Hiding}
\label{sec:props}

The properties are stated once at the vector-commitment level and specialise
to Merkle through $(\ROOutputSize,\MTSaltSize)$.  The first two are required;
the third is the optional, zero-knowledge case.

\begin{definition}[Completeness]\label{def:vc-complete}
$\VC$ is \emph{complete} if for every $\msgvec\in\MTAlphabet^{\MTMessageLength}$
and every $\idxset\subseteq[\MTMessageLength]$, with
$(\comm,\aux)\leftarrow\Commit(\msgvec)$ and
$\opening\leftarrow\Open(\aux,\idxset)$,
\[
  \Pr\!\left[\VCCheck(\comm,\idxset,\msgvec_{\idxset},\opening)=1\right]=1,
\]
where $\msgvec_{\idxset}$ is the restriction of $\msgvec$ to $\idxset$.  For
Merkle this is the honest commit--open--check round trip.
\end{definition}

\begin{definition}[Position binding]\label{def:vc-binding}
$\VC$ is \emph{position binding} if no efficient adversary $\adv$ outputs
$(\comm,\idxset_0,\idxset_1,\MTMessageSubVector_0,\MTMessageSubVector_1,
\opening_0,\opening_1)$ with both
$\VCCheck(\comm,\idxset_b,\MTMessageSubVector_b,\opening_b)=1$ yet
$\MTMessageSubVector_0[i]\ne\MTMessageSubVector_1[i]$ for some
$i\in\idxset_0\cap\idxset_1$, except with probability $\negl$.  The commitment
$\comm$ is adversarial: there is no honest-generation requirement.
\end{definition}

A verifier never inspects the whole vector; it opens a chosen subset $\idxset$
and checks only the returned entries.  Position binding guarantees each queried
answer is the unique value the commitment can be opened to, fixed the moment
$\comm$ is published and before the query set is chosen, so a soundness argument
may treat it as fixed; a single conflict at some $i\in\idxset_0\cap\idxset_1$ is
the local witness that binding has failed.

\begin{definition}[Hiding]\label{def:vc-hiding}
A randomised $\VC$ (whose $\Commit$ draws salts from a space $\rndspace$) is
\emph{hiding} if for any two messages $\msgvec_0,\msgvec_1$ of length
$\MTMessageLength$ and any adversary that, after seeing $\comm$, adaptively
requests openings of positions on which $\msgvec_0$ and $\msgvec_1$ agree, the
advantage in deciding which message was committed is at most $\negl$.
\end{definition}

Unlike binding, which is structural, hiding rests entirely on the salt space
$\rndspace$, which must be both \emph{well sampled} (drawn from a sound
generator and mixed into each leaf) and \emph{large enough} to resist search
within the security budget; either failure is fatal and neither repairs the
other.  For Merkle, the root names two refinements that \cref{sec:hiding}
discharges.

\begin{definition}[Root hiding]\label{def:root-hiding}
For a fixed $\msgvec$, the statistical distance between the real root (sampled
with fresh commitment randomness) and a simulated root produced without
$\msgvec$ is at most $\negl$.
\end{definition}

\begin{definition}[Selective-opening hiding]\label{def:selective-hiding}
After the commitment and the openings of a selected position set, the adversary
learns nothing about the unopened positions beyond what the disclosed values
already imply.
\end{definition}

Root hiding is the base step; selective-opening hiding is the stronger target,
because it must also hide through the disclosed salts and authentication paths.
A deterministic salt space satisfies neither---the receiver can recompute
candidate leaf commitments---so randomised salts are a prerequisite, not an
optional improvement.

\takeaway{The boundary fixes \emph{what} a vector commitment and its Merkle
instantiation guarantee: completeness, position binding at every queried
location, and optional hiding governed by the salt space.  \cref{sec:results}
turns each into an executable, Lean-checked contract and derives the parameters
that make it hold.}
